\documentclass[11pt]{amsart}

\usepackage[utf8]{inputenc}
\usepackage[T1]{fontenc}
\usepackage{amsmath,amssymb,amsthm}
\usepackage{hyperref}
\usepackage{xcolor}
\usepackage{listings}
\usepackage{booktabs}
\usepackage{enumitem}
\usepackage[margin=1in]{geometry}

% Lean 4 syntax highlighting
\lstdefinelanguage{Lean4}{
  morekeywords={def,theorem,lemma,structure,inductive,namespace,end,where,
    import,open,variable,instance,noncomputable,abbrev,example,
    by,intro,obtain,exact,simp,omega,cases,induction,refine,rw,
    have,let,show,change,apply,fun,match,with,if,then,else,do,
    return,sorry,Prop,Type,Sort,Bool,Nat,Fin,Set,true,false,
    extends,deriving},
  sensitive=true,
  morecomment=[l]{--},
  morecomment=[n]{/-}{-/},
  morestring=[b]",
  literate=
    {→}{{$\to$}}1
    {←}{{$\leftarrow$}}1
    {⟨}{{$\langle$}}1
    {⟩}{{$\rangle$}}1
    {∀}{{$\forall$}}1
    {∃}{{$\exists$}}1
    {¬}{{$\lnot$}}1
    {≤}{{$\leq$}}1
    {≥}{{$\geq$}}1
    {∧}{{$\land$}}1
    {∨}{{$\lor$}}1
    {×}{{$\times$}}1
    {ℕ}{{$\mathbb{N}$}}1
    {∘}{{$\circ$}}1
    {λ}{{$\lambda$}}1
    {α}{{$\alpha$}}1
    {β}{{$\beta$}}1
    {σ}{{$\sigma$}}1
    {₁}{{$_1$}}1
    {₂}{{$_2$}}1
}
\lstset{
  language=Lean4,
  basicstyle=\small\ttfamily,
  keywordstyle=\bfseries\color{blue!70!black},
  commentstyle=\itshape\color{green!50!black},
  stringstyle=\color{red!60!black},
  breaklines=true,
  columns=flexible,
  frame=single,
  framerule=0.4pt,
  xleftmargin=2em,
  xrightmargin=1em,
  aboveskip=0.8em,
  belowskip=0.8em,
  numbers=none,
  numberstyle=\tiny\color{gray},
  captionpos=b,
}

\theoremstyle{plain}
\newtheorem{theorem}{Theorem}[section]
\newtheorem{lemma}[theorem]{Lemma}
\newtheorem{proposition}[theorem]{Proposition}
\newtheorem{corollary}[theorem]{Corollary}

\theoremstyle{definition}
\newtheorem{definition}[theorem]{Definition}
\newtheorem{remark}[theorem]{Remark}

\newcommand{\BoolFn}{\mathsf{BoolFn}}
\newcommand{\Circ}{\mathsf{Circuit}}
\newcommand{\size}{\mathsf{size}}
\newcommand{\eval}{\mathsf{eval}}
\newcommand{\computes}{\mathsf{computes}}
\newcommand{\polyb}{\mathsf{PolyBounded}}
\newcommand{\polye}{\mathsf{PolyInExp}}
\newcommand{\Ppoly}{\mathsf{P/poly}}
\newcommand{\PRF}{\mathsf{PRF}}
\newcommand{\Lean}{\textsf{Lean\,4}}

\title[Natural Proofs Barrier in Lean 4]%
  {A Machine-Verified Natural Proofs Barrier in Lean~4}

\author{Alex Chengyu Li}
\address{Independent Researcher, Tokyo, Japan}
\email{cl7076@nyu.edu}

\begin{document}

\begin{abstract}
The natural proofs barrier (Razborov--Rudich 1997) has had no
machine-verified formalization in any proof assistant. We formalize it
in \Lean{}: if pseudorandom function generators exist, then no natural
combinatorial property is useful against $\Ppoly$. The formalization
is approximately 400 lines of \Lean{} with no unresolved proof
obligations.

The proof is short once definitions are chosen appropriately.
We identify definition choices that make the barrier proof mechanizable
without probabilistic infrastructure and analyze why simpler
alternatives fail. We also formalize the
Shannon counting argument: most Boolean functions require large circuits.
\end{abstract}

\maketitle

%% ===================================================================
\section{Introduction}
\label{sec:intro}

The natural proofs barrier, established by Razborov and Rudich~\cite{RR97},
is one of three fundamental obstacles to proving circuit lower bounds
(alongside relativization~\cite{BGS75} and algebrization~\cite{AW09}).
It shows that any proof technique satisfying two conditions,
\emph{constructivity} (efficiently checkable) and \emph{largeness}
(satisfied by many functions), cannot
establish superpolynomial circuit lower bounds, assuming pseudorandom
function generators exist. This barrier has shaped complexity theory for
nearly three decades, ruling out ``natural'' approaches to problems like
$\mathsf{P}$ vs.\ $\mathsf{NP}$.

No machine-verified formalization exists in any proof assistant:
not in Mathlib, Coq, Isabelle/AFP, or Agda (Section~\ref{sec:related}).

Our formalization targets \Lean{} with the Mathlib library. The barrier
theorem compiles to a six-line tactic proof. The difficulty is in
designing the definitions the proof operates on. We make three
contributions:
\begin{enumerate}[nosep]
\item The first machine-verified formalization of a complexity-theoretic
  barrier result.
\item A systematic account of definition choices: we document
  three definition choices that enable the mechanized proof and show why
  natural alternatives fail (Section~\ref{sec:design}).
\item A formalization of the Shannon counting argument, reusing the same
  circuit abstraction (Section~\ref{sec:shannon}).
\end{enumerate}

Section~\ref{sec:relationship} describes what is preserved from the
original proof and what is reformulated for mechanizability.


%% ===================================================================
\section{Informal Background}
\label{sec:background}

We briefly recall the statement and proof of the natural proofs
barrier, following Razborov and Rudich~\cite{RR97}.

\begin{definition}
A \emph{Boolean function} on $n$ variables is a map
$f \colon \{0,1\}^n \to \{0,1\}$.
A \emph{Boolean circuit} over the basis $\{\mathsf{AND}, \mathsf{OR},
\mathsf{NOT}\}$ computes such a function; its \emph{size} is the number
of gates.
\end{definition}

\begin{definition}
A \emph{combinatorial property} $\mathcal{P}$ of $n$-variable Boolean
functions is called \emph{natural} if it satisfies:
\begin{itemize}[nosep]
\item \textbf{Constructivity.} Whether $f$ satisfies $\mathcal{P}$ can be
  decided in time polynomial in $2^n$ (the truth table length).
\item \textbf{Largeness.} At least a $1/\mathrm{poly}(n)$ fraction of all
  $n$-variable Boolean functions satisfy $\mathcal{P}$.
\end{itemize}
The property is \emph{useful against} $\Ppoly$ if, for every polynomial
$p$, functions computable by circuits of size $\leq p(n)$ eventually
fail to satisfy $\mathcal{P}$.
\end{definition}

\begin{theorem}[Razborov--Rudich 1997]
\label{thm:informal}
If pseudorandom function generators exist, then no natural combinatorial
property is useful against $\Ppoly$.
\end{theorem}

\begin{proof}[Proof sketch]
Suppose $\mathcal{P}$ is natural and useful against $\Ppoly$. Let $G$ be
a PRF generator with polynomial circuit size bound $c \cdot n^d + c$.
Usefulness at this bound provides a cutoff $N$: for all $n \geq N$,
no function with circuits of size $\leq c \cdot n^d + c$ satisfies
$\mathcal{P}$. But constructivity and largeness together imply that
$\mathcal{P}$ would be an efficient distinguisher with non-negligible
advantage over the PRF, contradicting security. Concretely, at some
$n \geq N$, some PRF output $G(n, \mathit{seed})$ satisfies $\mathcal{P}$,
yet it has a circuit of size $\leq c \cdot n^d + c$. Contradiction.
\end{proof}


%% ===================================================================
\section{Formalization}
\label{sec:formalization}

Table~\ref{tab:files} summarizes the formalization structure.
All definitions and proofs compile under \Lean{} v4.29.0-rc6 with
Mathlib.

\begin{table}[ht]
\centering
\caption{Formalization structure.}
\label{tab:files}
\begin{tabular}{@{}lll@{}}
\toprule
File & Contents & Key definitions \\
\midrule
\texttt{Basic.lean} & Boolean functions & $\BoolFn$, \texttt{truthTable}, \texttt{card\_boolFn} \\
\texttt{Circuit.lean} & Circuits & $\Circ$, $\eval$, $\size$, $\computes$, \texttt{exists\_circuit} \\
\texttt{ComplexityClass.lean} & $\Ppoly$ & \texttt{CircuitFamily}, $\polyb$, $\polye$ \\
\texttt{NaturalProof.lean} & Natural properties & \texttt{CombProperty}, constructive, large, useful \\
\texttt{PRF.lean} & PRF hypothesis & \texttt{PRFGenerator}, \texttt{PRFHypothesis} \\
\texttt{NaturalProofBarrier.lean} & Main theorem & \texttt{natural\_proof\_barrier} \\
\texttt{Shannon.lean} & Counting bound & \texttt{finite\_bounded\_circuit}, \texttt{exists\_hard\_function} \\
\bottomrule
\end{tabular}
\end{table}

\subsection{Boolean functions and circuits}

A Boolean function on $n$ variables is a Lean function
$(\mathsf{Fin}\; n \to \mathsf{Bool}) \to \mathsf{Bool}$.
This representation matches circuit evaluation directly: no encoding
or decoding is needed at the interface between functions and circuits.

\begin{lstlisting}
abbrev BoolFn (n : ℕ) := (Fin n → Bool) → Bool
\end{lstlisting}

Circuits are an inductive type over the standard basis:

\begin{lstlisting}
inductive Circuit (n : ℕ) where
  | var   : Fin n → Circuit n
  | const : Bool → Circuit n
  | and   : Circuit n → Circuit n → Circuit n
  | or    : Circuit n → Circuit n → Circuit n
  | not   : Circuit n → Circuit n
\end{lstlisting}

Evaluation (\lstinline{Circuit.eval}) proceeds by structural recursion.
Gate count (\lstinline{Circuit.size}) counts AND, OR, and NOT nodes;
inputs and constants contribute zero. A circuit $c$ \emph{computes}
$f$ when $\forall x,\; c.\eval\; x = f\; x$.

\subsection{Complexity classes}

A circuit family assigns one circuit per input length.
The size function is \emph{polynomially bounded} if
$\exists\, c\, d,\; \forall n,\; s(n) \leq c \cdot n^d + c$.
A function sequence is in $\Ppoly$ when computed by a polynomially
bounded circuit family.

For constructivity, we need a separate notion: the checker circuit
size is polynomial in the \emph{truth table length} $2^n$, not in
the number of variables $n$. We formalize this as $\polye$:

\begin{lstlisting}
def PolyInExp (s : ℕ → ℕ) : Prop :=
  ∃ c d : ℕ, ∀ n, s n ≤ c * (2 ^ n) ^ d + c
\end{lstlisting}

The difference between $\polyb$ and $\polye$ is exponential and must
be enforced at the type level (Section~\ref{sec:polyinexp}).

\subsection{Natural properties}

A combinatorial property bundles a predicate with decidability:

\begin{lstlisting}
structure CombProperty where
  prop : (n : ℕ) → BoolFn n → Prop
  decidable : (n : ℕ) → DecidablePred (prop n)
\end{lstlisting}

Constructivity requires a checker circuit family of $\polye$ size
that agrees with the property on truth table inputs.
Largeness requires that the fraction of satisfying functions is
$\geq 1/q(n)$ for some polynomially bounded $q$.

Usefulness is formalized \emph{per polynomial bound}, not per
function family (see Section~\ref{sec:family-level} for why):

\begin{lstlisting}
def CombProperty.usefulAgainstPPoly (P : CombProperty) : Prop :=
  ∀ c d : ℕ, ∃ N, ∀ n, n ≥ N →
    ∀ (f : BoolFn n) (circ : Circuit n),
      circ.computes f → circ.size ≤ c * n ^ d + c →
        ¬P.prop n f
\end{lstlisting}

A \texttt{NaturalProperty} extends \texttt{CombProperty} with proofs
of constructivity and largeness.

\subsection{PRF hypothesis}

A \texttt{PRFGenerator} maps seeds to Boolean functions via
polynomial-size circuits, with a \emph{uniform} bound
$c \cdot n^d + c$ over all seeds:

\begin{lstlisting}
structure PRFGenerator where
  seedLen  : ℕ → ℕ
  generate : (n : ℕ) → (Fin (seedLen n) → Bool) → BoolFn n
  polySize : ∃ c d : ℕ, ∀ n (seed : Fin (seedLen n) → Bool),
    ∃ circ : Circuit n,
      circ.computes (generate n seed) ∧
      circ.size ≤ c * n ^ d + c
\end{lstlisting}

The PRF hypothesis states the consequence of PRF security used in the
barrier proof (see Section~\ref{sec:prob-prf} for justification):

\begin{lstlisting}
def PRFHypothesis : Prop :=
  ∃ G : PRFGenerator,
    ∀ (P : CombProperty),
      P.constructive → P.large →
      ∀ N : ℕ, ∃ n, n ≥ N ∧
        ∃ seed : Fin (G.seedLen n) → Bool,
          P.prop n (G.generate n seed)
\end{lstlisting}

\subsection{The barrier theorem}

The main theorem and its complete proof:

\begin{lstlisting}
theorem natural_proof_barrier (h : PRFHypothesis) :
    ∀ P : NaturalProperty,
      ¬P.toCombProperty.usefulAgainstPPoly := by
  intro P hU
  obtain ⟨G, hG⟩ := h
  obtain ⟨c_G, d_G, heff⟩ := G.polySize
  obtain ⟨N, hN⟩ := hU c_G d_G
  obtain ⟨n, hn, seed, hsat⟩ :=
    hG P.toCombProperty P.isConstructive P.isLarge N
  obtain ⟨circ, hcomp, hsize⟩ := heff n seed
  exact hN n hn (G.generate n seed) circ hcomp hsize hsat
\end{lstlisting}

The proof proceeds in six steps:
\begin{enumerate}[nosep]
\item Extract a PRF generator $G$ from the hypothesis.
\item Obtain its uniform circuit size bound $(c_G, d_G)$.
\item Apply usefulness at $(c_G, d_G)$ to get a cutoff $N$.
\item The PRF hypothesis, applied to the constructive and large property,
  yields $n \geq N$ and a seed such that $G(n, \mathit{seed})$ satisfies $\mathcal{P}$.
\item The efficiency bound gives a circuit for $G(n, \mathit{seed})$ of
  size $\leq c_G \cdot n^{d_G} + c_G$.
\item Usefulness at $n \geq N$ says no function with such a small circuit
  satisfies $\mathcal{P}$, contradicting step~4.
\end{enumerate}


%% ===================================================================
\section{Design Space: Rejected Alternatives}
\label{sec:design}

Each definition in Section~\ref{sec:formalization} avoids a
mechanization obstacle that blocks simpler formulations.

\subsection{Family-level usefulness}
\label{sec:family-level}

A simpler formulation of usefulness quantifies over function
families in $\Ppoly$:
\[
  \forall f \in \Ppoly,\; \exists N,\; \forall n \geq N,\;
  \neg\mathcal{P}(f_n).
\]
This formulation fails mechanically: the PRF generator produces outputs indexed
by \emph{seed}, so step~4 of the proof yields a specific seed
$s_0$ at a specific $n_0$. But usefulness provides a cutoff $N$
that depends on the \emph{family} $f$. Different seeds yield
different families, each with its own cutoff $N_s$. To reach a
contradiction, we need a uniform cutoff that works for all seeds
simultaneously, but over infinitely many seeds, no finite maximum
exists.

\paragraph{Fix.}
Quantify over polynomial bounds $(c, d)$ directly:
\[
  \forall c\, d,\; \exists N,\; \forall n \geq N,\;
  \bigl(\size \leq c \cdot n^d + c \implies \neg\mathcal{P}\bigr).
\]
The PRF generator's \texttt{polySize} field pins a single $(c_G, d_G)$.
This yields one cutoff $N$ for all seeds, because all seeds share the
same circuit size polynomial. The contradiction closes uniformly.

\subsection{Probabilistic PRF definition}
\label{sec:prob-prf}

The textbook PRF definition uses $\varepsilon$-indistinguishability:
\[
  \bigl|\Pr[D(G(k)) = 1] - \Pr[D(f) = 1]\bigr| \leq \varepsilon(n),
\]
where $f$ is drawn uniformly from all $n$-variable Boolean functions.
Formalizing this requires:
\begin{itemize}[nosep]
\item Probability distributions over $\BoolFn\; n$ (uniform measure on
  a space of size $2^{2^n}$).
\item Negligible functions (asymptotic rate conditions).
\item The distinguisher-to-property reduction: if $\mathcal{P}$ is
  constructive, large, and rejects all PRF outputs, one can build
  a distinguisher with non-negligible advantage.
\end{itemize}
All three require Mathlib infrastructure that does not yet exist.
The distinguisher reduction is the only step that uses the PRF
definition's \emph{internal structure}. The barrier proof uses only
the \emph{consequence}: if $\mathcal{P}$ is constructive and large, it
cannot reject all PRF outputs at arbitrarily large~$n$.

\paragraph{Fix.}
State this consequence directly as \texttt{PRFHypothesis}. The barrier
proof goes through in six tactic steps. The probabilistic infrastructure
would add hundreds of lines of measure-theoretic scaffolding that the
barrier theorem does not use.

The consequence form is implied by standard PRF security, so no logical
content is lost. Any formalization starting from the standard definition
must derive this consequence as an intermediate lemma.

\subsection{PolyInExp vs.\ PolyBounded}
\label{sec:polyinexp}

Constructivity means the checker runs in time polynomial in the
\emph{truth table length} $2^n$, not in the number of variables $n$.
The informal proof says ``polynomial time'' without specifying
the parameter. Using $\polyb$ instead of $\polye$ for the checker
makes the constructivity condition exponentially too weak: the
checker would need to run in time $\mathrm{poly}(n)$, which is
sublinear in its $2^n$-bit input.

In the formalization, this distinction is enforced at the type level:
\texttt{constructive} uses $\polye$ while \texttt{usefulAgainstPPoly}
uses $\polyb$. A type error would result from swapping them.


%% ===================================================================
\section{Relationship to Razborov--Rudich (1997)}
\label{sec:relationship}

This section distinguishes faithful translations from machine-oriented
reformulations in the formalization.

\paragraph{Faithful translations.}
The contradiction structure is identical to the original: usefulness
provides a cutoff below which no small-circuit function satisfies
$\mathcal{P}$; the PRF hypothesis produces a small-circuit function
satisfying $\mathcal{P}$ above that cutoff. Boolean functions,
circuits, and the AND/OR/NOT basis match the original. Constructivity
and largeness retain their combinatorial content.

\paragraph{Reformulations.}
The PRF hypothesis is stated in consequence form (Section~\ref{sec:prob-prf})
and usefulness per polynomial bound (Section~\ref{sec:family-level}).
Both reformulations drop technical apparatus (probability distributions,
circuit family indirection) that the barrier proof does not use,
while preserving the logical content of each step.

\paragraph{Scope of reformulation.}
The overall statement is faithful to the Razborov--Rudich barrier theorem.
The reformulations concern how the proof \emph{accesses} its hypotheses.


%% ===================================================================
\section{Shannon Counting Bound}
\label{sec:shannon}

We also formalize the Shannon counting argument: most Boolean functions
require large circuits.
This result reuses the $\Circ$, $\size$, $\BoolFn$, and $\computes$
definitions from the barrier formalization without modification.

\begin{theorem}[\texttt{finite\_bounded\_circuit}]
For all $n, s$, the set $\{c : \Circ\; n \mid c.\size \leq s\}$
is finite.
\end{theorem}

By induction on $s$. In the base case ($s = 0$),
circuits of size zero are exactly variables and constants, a finite
set. In the inductive step, circuits of size $\leq s+1$ are either
in the inductive hypothesis set or are built from subcircuits of
size $\leq s$ via AND, OR, or NOT, each a finite image of a
finite product.

\begin{theorem}[\texttt{finite\_computable\_bounded}]
For all $n, s$, the set $\{f : \BoolFn\; n \mid \exists\, c,\;
c.\computes\; f \land c.\size \leq s\}$ is finite.
\end{theorem}

This follows immediately: the set is a subset of the image of
bounded circuits under evaluation.

\begin{theorem}[\texttt{exists\_hard\_function}]
If the number of functions computable by circuits of size $\leq s$
is less than the total number of $n$-variable Boolean functions
($= 2^{2^n}$), then there exists a function requiring circuits
of size ${>}\, s$.
\end{theorem}

Suppose for contradiction that every function has a circuit of
size $\leq s$. Then $\mathsf{Finset.univ}$ (all $2^{2^n}$ functions)
is a subset of the finite computable set, contradicting
the cardinality hypothesis.


%% ===================================================================
\section{Related Work}
\label{sec:related}

\paragraph{Complexity theory in proof assistants.}
The closest precedent is the formalization of the Cook--Levin theorem
(NP-completeness of SAT) in Coq by Forster and
Kunze~\cite{ForsterKunze21}, which used the call-by-value
$\lambda$-calculus L as a computational model rather than Turing
machines. Balbach~\cite{Balbach23} independently formalized Cook--Levin
in Isabelle/HOL using Turing machines. Forster, Kunze, and
Lauermann~\cite{ForsterKL22} formalized synthetic Kolmogorov
complexity in Coq. Those results are positive theorems:
NP-completeness, existence of incompressible strings.
No complexity barrier has been formalized previously.

\paragraph{Barrier results.}
The relativization barrier~\cite{BGS75} and the
algebrization barrier~\cite{AW09} also remain unformalized.
Relativization would require formalizing oracle computation, which
introduces new types and abstraction layers; we discuss this as
future work.

\paragraph{Cryptography in proof assistants.}
SSProve~\cite{SSProve23} provides a framework for machine-checked
cryptographic proofs in Coq, including PRF-based encryption, using
the full probabilistic definition. Our consequence-form PRF hypothesis
sidesteps this overhead, since the barrier proof uses only the
consequence, as discussed in Section~\ref{sec:prob-prf}.


%% ===================================================================
\section{Conclusion}
\label{sec:conclusion}

This paper formalizes the natural proofs barrier in \Lean{},
with no \texttt{sorry} obligations and approximately 400 lines of
code. The definitions do more work than the proofs: choosing interfaces
that avoid probabilistic infrastructure is the main obstacle in
mechanizing this class of results.

The design space documented in Section~\ref{sec:design} records which
definition patterns support mechanized proofs of barrier results and
which fail.

\paragraph{Future directions.}
The Shannon counting bound (Section~\ref{sec:shannon}) shows that the
circuit abstraction applies beyond the barrier theorem. Two extensions
that build on the current definitions without modifying them are:
\begin{itemize}[nosep]
\item \textbf{Circuit size bounds.} Concrete upper and lower bounds
  on the number of circuits of given size, enabling explicit
  Shannon-type lower bounds for specific parameters.
\item \textbf{Karp--Lipton collapse.} If SAT has polynomial-size
  circuits, then the polynomial hierarchy collapses.
\end{itemize}
The relativization and algebrization barriers require oracle
computation types, and thus new core abstractions rather than
additional lemmas.

\paragraph{Artifact.}
The complete \Lean{} source code is available at
\url{https://github.com/crabsatellite/lean4-natural-proofs-barrier}.
This paper is archived at Zenodo (DOI: \texttt{10.5281/zenodo.19132831}).


\begin{thebibliography}{99}

\bibitem{AW09}
S.~Aaronson and A.~Wigderson.
\newblock Algebrization: a new barrier in complexity theory.
\newblock \emph{ACM Transactions on Computation Theory}, 1(1):2:1--2:54, 2009.

\bibitem{Balbach23}
F.~Balbach.
\newblock The {Cook--Levin} theorem in {Isabelle/HOL}.
\newblock In \emph{Archive of Formal Proofs}, 2023.

\bibitem{BGS75}
T.~Baker, J.~Gill, and R.~Solovay.
\newblock Relativizations of the {P} $=?$ {NP} question.
\newblock \emph{SIAM Journal on Computing}, 4(4):431--442, 1975.

\bibitem{ForsterKunze21}
Y.~Forster and F.~Kunze.
\newblock Mechanising complexity theory: the {Cook--Levin} theorem in {Coq}.
\newblock In \emph{ITP 2021}, volume 193 of \emph{LIPIcs}, pages 20:1--20:18, 2021.

\bibitem{ForsterKL22}
Y.~Forster, F.~Kunze, and N.~Lauermann.
\newblock Synthetic {Kolmogorov} complexity in {Coq}.
\newblock In \emph{ITP 2022}, volume 237 of \emph{LIPIcs}, pages 12:1--12:19, 2022.

\bibitem{RR97}
A.~A. Razborov and S.~Rudich.
\newblock Natural proofs.
\newblock \emph{Journal of Computer and System Sciences}, 55(1):24--35, 1997.

\bibitem{SSProve23}
P.~Barbosa, G.~Barthe, B.~Gr{\'e}goire, A.~Koutsos, and P.-Y.~Strub.
\newblock {SSProve}: a foundational framework for modular cryptographic proofs
  in {Coq}.
\newblock \emph{ACM Transactions on Programming Languages and Systems},
  45(1):5:1--5:44, 2023.

\end{thebibliography}

\end{document}
